% Generator: Claude Opus 5 (High)
% Section added 2026-08-11. Both exercises are transcribed from the Probeprfg mock exams (see the
% % Source: comments below); the definition, the connecting prose and the remarks are written for
% this document, which had no notion of a centroid anywhere. The two problems come from different
% papers and are placed together because each is the other's payoff: one computes the centroid,
% the other says what the centroid is for.

\section{Centroids}
\label{sec:centroids}

% Transition: Claude Opus 5 (High)
Every integral in this chapter so far has been computed for its own sake. The simplest thing one
actually does with an integral over a region is average a coordinate over it, and the resulting
point is worth a name.

\begin{definition}[Centroid of a plane region]
\label{def:centroid}
Let $B \subseteq \mathbb{R}^2$ be Jordan measurable with $\vol_2(B) > 0$. The \newterm{centroid}
\germanfor{Fl\"achenschwerpunkt} of $B$ is the point $z_B := (x_B, y_B)$ with
\[x_B := \frac{1}{\vol_2(B)}\int_B x \,d\vol_2(x,y), \qquad
  y_B := \frac{1}{\vol_2(B)}\int_B y \,d\vol_2(x,y).\]
\end{definition}

The same formula defines the centroid of a Jordan measurable set in any $\mathbb{R}^n$, one
coordinate at a time; nothing below uses the restriction to the plane except the triangle.

\begin{remark}[Why the normalisation is there]
Without the factor $1/\vol_2(B)$ the two integrals would scale with the size of $B$ as well as
with its position, and doubling a region would move its \qt{centre}. Dividing by the area makes
$z_B$ an \emph{average} of the coordinate functions rather than a total, and the immediate
consequence is that $z_B$ is where one expects it to be: for a region symmetric about a line, the
centroid lies on that line, because the coordinate perpendicular to it integrates to zero.
\end{remark}

\begin{center}
% Generator: Claude Opus 5 (High) --- FIG-20-02: the two halves of the exercise below, drawn.
% This section defines a *point* and then asserts two things about where it sits, with no picture
% anywhere: the reader meets the weighted-average formula of part (a) as an identity between two
% integrals and never sees that it puts $z_B$ on the segment joining the two sub-centroids.
% Both panels are exact, not schematic --- the coordinates below are the computed centroids, so
% the picture can be measured against the formulas it illustrates.
% `lbl' backs a label with a soft white plate so it stays readable where it lands on a shaded
% region or a box edge -- the same idiom as the surface figure in ch. 21, and the thing that
% decides whether these read on a grey e-reader screen rather than only at 300dpi in colour.
\begin{tikzpicture}[>=Latex, scale=1.0,
  lbl/.style={fill=white, fill opacity=0.8, text opacity=1, inner sep=1.2pt,
              rounded corners=1pt}]

  % ---------- Left: centroids add, as an area-weighted average ----------
  % B_1 = [0,2]x[0,1.8], area 3.6, centroid (1, 0.9)
  % B_2 = [2,3]x[0,0.9], area 0.9, centroid (2.5, 0.45)
  % weight of B_2 is 0.9/4.5 = 1/5 exactly, so
  % z_B = z_1 + (1/5)(z_2 - z_1) = (1.3, 0.81)
  \node[font=\small\bfseries, text=TextBoldColor] at (1.5, 3.0) {Centroids add};

  \fill[blue!10, opacity=0.7] (0,0) rectangle (2,1.8);
  \fill[orange!12, opacity=0.7] (2,0) rectangle (3,0.9);
  \draw[blue!75!black, thick] (0,0) rectangle (2,1.8);
  \draw[orange!90!black, thick] (2,0) rectangle (3,0.9);
  \node[blue!80!black, font=\scriptsize] at (0.45, 1.5) {$B_1$};
  \node[orange!90!black, font=\scriptsize] at (2.75, 0.68) {$B_2$};

  % the segment joining the two sub-centroids, and z_B on it
  \draw[gray!70!black, dashed] (1,0.9) -- (2.5,0.45);
  \fill[blue!75!black] (1,0.9) circle (1.5pt)
    node[lbl, above left=-2pt, font=\scriptsize] {$z_{B_1}$};
  \fill[orange!90!black] (2.5,0.45) circle (1.5pt)
    node[lbl, below right=-2pt, font=\scriptsize] {$z_{B_2}$};
  \fill[red!80!black] (1.3,0.81) circle (1.8pt);
  % text=, not a bare colour, which would override the fill=white in `lbl'.
  \node[lbl, text=red!80!black, font=\scriptsize, above right=-2pt] at (1.3,0.81) {$z_B$};

  \node[font=\scriptsize, align=center] at (1.5,-0.75)
    {$\vol_2(B_2)/\vol_2(B) = \tfrac15$, and $z_B$ sits\\one fifth of the way along};

  % ---------- Right: the centroid of a triangle ----------
  % o=(0,0), p=(3,0), q=(1.2,2.4); centroid = (4.2/3, 2.4/3) = (1.4, 0.8)
  \begin{scope}[xshift=5.0cm]
    \node[font=\small\bfseries, text=TextBoldColor] at (1.5, 3.0) {\dots\ and a triangle's};

    \fill[blue!10, opacity=0.7] (0,0) -- (3,0) -- (1.2,2.4) -- cycle;
    \draw[blue!75!black, thick] (0,0) -- (3,0) -- (1.2,2.4) -- cycle;

    % the three medians, which meet at the centroid
    \draw[gray!70!black, thin] (0,0) -- (2.1,1.2);
    \draw[gray!70!black, thin] (3,0) -- (0.6,1.2);
    \draw[gray!70!black, thin] (1.2,2.4) -- (1.5,0);

    % midpoints
    \fill[gray!70!black] (2.1,1.2) circle (1.0pt);
    \fill[gray!70!black] (0.6,1.2) circle (1.0pt);
    \fill[gray!70!black] (1.5,0) circle (1.0pt);

    \fill[black] (0,0) circle (1.5pt) node[below left=-2pt, font=\scriptsize] {$o$};
    \fill[black] (3,0) circle (1.5pt) node[below right=-2pt, font=\scriptsize] {$p$};
    \fill[black] (1.2,2.4) circle (1.5pt) node[above=1pt, font=\scriptsize] {$q$};

    \fill[red!80!black] (1.4,0.8) circle (1.8pt);
    \node[lbl, text=red!80!black, font=\scriptsize, right=1pt] at (1.4,0.8) {$z_B$};

    \node[font=\scriptsize, align=center] at (1.5,-0.75)
      {the three medians meet at\\$z_B = \tfrac13(o+p+q)$};
  \end{scope}

\end{tikzpicture}
\end{center}

% Source: old_exams/fs2023/Probeprfg3.pdf (Probeprüfung Analysis II, 31 January 2023),
% Aufgabe 11, p. 4
% Extractor: Claude Opus 5 (High)
% The paper names no lecturer; see old-exam-mining.md.
\begin{exercise}[Centroids add, and the centroid of a triangle]
\label{ex:centroid_additivity_and_triangle}
Let $B \subseteq \mathbb{R}^2$ be Jordan measurable with $\vol_2(B) > 0$.
\begin{enumerate}[label=\textbf{(\alph*)}]
  \item Suppose $B = \bigcup_{i=1}^n B_i$ for pairwise disjoint Jordan measurable sets
    $B_1,\dots,B_n \subseteq \mathbb{R}^2$, each with $\vol_2(B_i) > 0$. Show that
    \[z_B = \sum_{i=1}^n \frac{\vol_2(B_i)}{\vol_2(B)}\,z_{B_i}.\]
  \item Show that if $B$ is a triangle with vertices $o = (0,0)$, $p = (r,0)$ and $q = (s,t)$,
    where $r,s,t > 0$, then $z_B = \tfrac13(o+p+q)$.
\end{enumerate}
\exinfo{This exercise is taken from the mock examination \emph{Probepr\"ufung Analysis II} of
31 January 2023, Problem 11. That paper names no lecturer.}
\exsol{sol:centroid_additivity_and_triangle}
\end{exercise}

% Source: old_exams/fs2023/Probeprfg2.pdf (Probeprüfung Analysis II, 17 August 2022, revised),
% Aufgabe 1, p. 1
% Extractor: Claude Opus 5 (High)
% The paper names no lecturer; see old-exam-mining.md. The original asks only for the minimising
% pair; the two-part split, and part (b) in particular, are added here, because the general
% statement is what makes the computation worth doing and the exam's own solution differentiates
% under the integral sign without ever saying that the answer is the centroid.
\begin{exercise}[What the centroid minimises]
\label{ex:centroid_minimises_second_moment}
Let $B := \{(x,y) \in \mathbb{R}^2 \mid x \geq 0,\ y \geq 0,\ x^2+y^2 \leq 1\}$ be the closed
quarter disc, and for $(s,t) \in \mathbb{R}^2$ put
\[I(s,t) := \int_B (x-s)^2 + (y-t)^2 \,d\vol_2(x,y).\]
\begin{enumerate}[label=\textbf{(\alph*)}]
  \item Determine the pair $(s,t)$ at which $I$ is minimal.
  \item Show that the answer to part \textbf{(a)} is no accident: for \emph{any} Jordan
    measurable $B \subseteq \mathbb{R}^2$ with $\vol_2(B) > 0$, the function $I$ is minimised
    exactly at $(s,t) = z_B$, and at no other point.
\end{enumerate}
\exinfo{Part \textbf{(a)} is taken from the mock examination \emph{Probepr\"ufung Analysis II} of
17 August 2022 (revised), Problem 1; that paper names no lecturer. Part \textbf{(b)} is added
here, since the general statement is what the computation is really about.}
\exsol{sol:centroid_minimises_second_moment}
\end{exercise}
